% ====================================================================
% graphite_ica_chapter1.tex  —  Chapter 1 (v4, synthesis)
% Graphite anode ICA (dQ/dV) tail: effective-barrier + relaxation-length
%   spectrum kernel-integral theory, closed by Refs 6/7 propagator method.
% Date: 2026-05-28   Author: Project_Anode_Fit (Claude 측)
% --------------------------------------------------------------------
% Synthesis of (a) Codex 강점: relaxation-length spectrum kernel integral
%   (stretched tail), (b) Claude rigor: grounding tier (Assumption Ledger),
%   사용자 PhD Refs 6/7 closure, falsification + χ-discriminator,
%   charge-balance implicit V_n, Marcus bound.
% Governance: Charter v3 / MASTER_ROADMAP_v3 §9bis / ASSUMPTION_LEDGER_v3 (AL-#)
% Rule: Korean prose + English academic terms. smooth only (no 0->1 step jump).
% ====================================================================
\documentclass[11pt,a4paper]{article}
\usepackage[margin=22mm]{geometry}
\usepackage{kotex}
\usepackage{amsmath,amssymb,mathtools,bm}
\usepackage{booktabs,longtable,array,tabularx}
\usepackage{enumitem}
\usepackage{hyperref}
\usepackage{amsthm}
\usepackage{mdframed}
\usepackage{fancyhdr}
\usepackage{lastpage}
\usepackage{setspace}
\setstretch{1.13}
\setlength{\parskip}{0.4em}
\setlength{\parindent}{0pt}
\setlength{\headheight}{14pt}
\hypersetup{colorlinks=true, linkcolor=blue, urlcolor=blue, citecolor=blue,
  pdftitle={Graphite ICA tail: barrier-spectrum kernel-integral theory (Ch.1 v4)}}
\pagestyle{fancy}\fancyhf{}
\lhead{Graphite ICA tail — Ch.1 (v4 synthesis)}\rhead{\thepage/\pageref{LastPage}}
\renewcommand{\headrulewidth}{0.4pt}

\newtheorem*{groundingbox}{Grounding}
\newtheorem*{boundbox}{Validity bound}
\newtheorem*{flagbox}{Flagged (미확정)}
\newtheorem*{uservb}{사용자 verbatim}

\newcommand{\dd}{\mathrm{d}}
\newcommand{\eff}{\mathrm{eff}}
\newcommand{\eq}{\mathrm{eq}}
\newcommand{\app}{\mathrm{app}}
\newcommand{\dimedrive}{\mathrm{drive}}
\newcommand{\tot}{\mathrm{tot}}
\newcommand{\bg}{\mathrm{bg}}
\newcommand{\cell}{\mathrm{cell}}
\newcommand{\ext}{\mathrm{ext}}
\newcommand{\simplecase}{\mathrm{simple}}
\newcommand{\AL}[1]{\textsf{[AL-#1]}}

\title{\textbf{리튬이온전지 graphite anode 의 ICA(\,$dQ/dV$\,) tail 거동:\\
effective barrier 와 relaxation-length spectrum kernel-integral 이론}\\[0.4em]
\large Chapter 1 (v4 synthesis, 2026-05-28)}
\author{Project\_Anode\_Fit}
\date{2026-05-28}

\begin{document}
\maketitle

\begin{center}\begin{minipage}{0.93\textwidth}
\small \textbf{Abstract.} 리튬이온전지 graphite anode 의 incremental capacity analysis
(ICA), $dQ/dV$ 곡선에서 staging transition peak 뒤의 \emph{post-peak tail} 은 temperature
가 낮을수록 길게 늘어지고(stretched) 높을수록 짧게 끝난다. 본 Chapter 1 은 이 tail 을
\textbf{단일 relaxation length 가 아니라 relaxation-length spectrum 의 kernel integral}
로 해석한다. 출발은 \textbf{activation free-energy barrier 가 state 를 $0\!\to\!1$ 로
점프시키는 threshold 가 아니라 rate constant 를 정하는 hidden variable} 이라는 점이다
(step-function completion 금지). 각 local mode 는 effective barrier $\Delta G_\eff =
\Delta G_a - \chi\,\mathcal A$ (Brønsted--Evans--Polanyi / Butler--Volmer; Marcus 로 유효
범위 한정) 가 정하는 Eyring rate 로 equilibrium target 을 향해 relaxation 하며, 하나의
exponential kernel ($L=v_Q/k$) 을 만든다. 관측 tail 은 barrier distribution $\rho_G$ 가
\emph{지수적 barrier$\to$length 매핑}을 통해 만드는 relaxation-length spectrum $A_L$ 에
대한 kernel integral 이다. 이 self-consistent spectrum integral 의 closed-form 은
\textbf{사용자 박사학위 연구 (Refs~\cite{lee2011,son2013}) 의 propagator / ratio-
substitution 기법}으로 닫는다. low temperature 는 spectrum 을 large-$L$ 로, electrode-
potential assistance 는 short-$L$ 로 이동시켜 관측 tail 길이를 설명한다. equilibrium
target 의 함수 형태는 특정하지 않고(보수적; grounded lattice-gas 를 특수해로만), $V_n$ 은
외부 lookup 이 아니라 charge conservation 으로 implicit 결정한다. 모든 assumption 은
Assumption Ledger (AL-\#) 에 grounding 되며, 근거 없는 항목은 flagged 로만 둔다. tail 을
effective barrier 에 귀속하기에 앞서 polarization·transport·heterogeneous width 등 competing
source 를 구별하는 falsification protocol 과, barrier distribution 의 \emph{non-unique}
역산을 피하는 forward-prediction 원칙을 제시한다.
\end{minipage}\end{center}

\tableofcontents
\newpage

% ====================================================================
\section{Motivation 과 deliverable}\label{sec:intro}
% ====================================================================

\subsection{사용자 의도 (verbatim)}\label{sec:verbatim}
\begin{uservb}
``LIB의 ICA 분석에서 온도와 현시점의 전위 상태에 따라서 그 그래프의 개형이 특히
피크부분이 상변이에 의해서 나타나는데 그 상변이의 의해 나타나는 피크 꼬리의 개형이
늘어지느냐 빨리 떨어지느냐가 달라지는것을 확인했다. 특히 온도가 낮을때는 꼬리가
길고, 온도가 높으면 그 꼬리가 상대적으로 짧게 끝나는 현상이 관측되었어. 원래라면
그렇게 온도에 대해 약간 가우시안 피크 형상이 나타나고 더이상 진행이 안될 것으로
생각되는데 그 온도에 의한 배리어 외에 추가적인 극판 자체 전위에 따른 배리어를
낮추는 효과가 있다고 봤고 그 유효 배리어의 논리를 확정하고 그 것을 이용해 피크의
거동을 해석하는 피팅하는데 쓸수 있는 논리 식을 구하는게 이 작업의 핵심이야.''
\end{uservb}

\textbf{해석 정정 (사용자 2026-05-28).} (1) ``가우시안''은 \emph{sharp peak} 를 가리키는
예시이며 equilibrium 형태를 Gaussian 으로 확정한 것이 아니다 → 형태는 grounding 과 data
로 결정. (2) ``activation energy barrier''는 본 이론의 핵심으로 \emph{유지}하되, ``특정
지점을 넘으면 $0\!\to\!1$ 로 갑자기 점프''하는 step-function 식 \emph{비약은 금지} (모든
transition 은 smooth). 본 chapter 는 이 두 정정을 절대 기준으로 한다.

\subsection{관측 사실 (observation)}\label{sec:observation}
\begin{enumerate}[label=\textbf{(O\arabic*)}]
\item \textbf{Peak 출현.} graphite staging transition 영역에서 $dQ/dV_{n,\app}$ 가 sharp
  peak 를 갖는다 (staging sequence 가 $N_p\approx3\sim4$ effective transition 으로 fusion).
\item \textbf{Peak tail 의 stretched temperature 의존.} 각 peak 의 tail 이 $T$ 가
  낮을수록 길게 늘어지고(흔히 단일 exponential 보다 \emph{stretched}) 높을수록 짧게
  끝난다. \emph{단일 equilibrium peak width 만으로는} 이 $T$/potential 의존 tail 을
  설명하기 어렵다.
\end{enumerate}
본 Chapter 의 1차 초점은 (O2). peak area 정합은 fitting 단계로 이관 (P1).

\subsection{가설 (hypothesis)}\label{sec:hypothesis}
\textbf{(H1)} equilibrium target $\theta_{\eq,j}(V_n,T)$ 는 $V_n$ 에 대해 smooth 하게
변한다 ($0\!\to\!1$ 급점프 아님; \S\ref{sec:equilibrium}).

\textbf{(H2)} 각 local mode 의 진행 속도는 activation barrier $\Delta G_{a}$ 를 통과하는
thermally activated 과정이다 (\AL{1}).

\textbf{(H3, ★)} electrode potential $V_{n,\dimedrive}$ 가 그 barrier 를 $\chi\,\mathcal A$
만큼 낮춘다: $\Delta G_\eff = \Delta G_a - \chi\mathcal A$, $\mathcal A=s_\phi F(V_{n,\dimedrive}-U)$
(\AL{2}; 정식 \S\ref{sec:barrier}).

\textbf{(H4)} barrier 는 rate constant $k=\nu\exp(-\Delta G_\eff/RT)$ 를 정하고 (\AL{1}),
$\theta_j$ 는 target $\theta_{\eq,j}$ 를 finite rate 로 따라가 lag 를 만든다 (\AL{5}).
하나의 mode 는 length $L=v_Q/k$ 의 \emph{단일 exponential kernel} 을 만든다.

\textbf{(H5, ★ 핵심)} 실제 electrode 는 단일 barrier 가 아니라 \emph{barrier
distribution} $\rho_G$ 를 가진다. barrier$\to$length 의 \emph{지수} 매핑이 그 분산을
크게 확대하여 \emph{relaxation-length spectrum} $A_L$ 을 만들고, 관측 tail 은 그 spectrum
의 \textbf{kernel integral} 이다 (\S\ref{sec:spectrum}--\ref{sec:kernel}). 따라서 관측
tail 은 일반적으로 stretched 이다.

\textbf{(H6)} $T$ 가 낮으면 spectrum 이 large-$L$ 로 이동해 tail 이 길어지고, $T$ 가
높거나 potential assistance 가 크면 short-$L$ 로 이동해 tail 이 짧아진다 (O2; \S\ref{sec:tailT}).

\subsection{Deliverable}\label{sec:deliverable}
(H1)$\sim$(H6) 을 grounded derivation 으로 확정하고, ICA tail 의 $T$/potential 의존을
설명하는 \emph{fitting 가능한 closed-form 논리식}(spectrum kernel integral, Refs 6/7 로
closure)을 도출한다. 단 (i) 모든 assumption 은 AL-\# grounding, (ii) FLAGGED 는 established
로 쓰지 않음, (iii) tail 의 barrier 귀속 전에 competing source 의 falsification protocol
제시 + barrier distribution 의 non-unique 역산 회피, (iv) fitting solver 구현은 범위 외 (P1).

\subsection{사용자 박사 연구의 위치 (Refs 6/7) = kernel integral 의 closure}\label{sec:phd}
spectrum kernel integral 과 self-consistent relaxation 은 미지 함수가 적분 안팎에 동시에
나타나는 integral equation 을 만든다. 사용자 박사 연구 \cite{lee2011,son2013} 의
\emph{propagator / ratio-substitution} 기법은 이런 적분방정식을 closed-form 으로 닫는
도구이며 (\S\ref{sec:closure}), 본 Chapter 의 load-bearing 수단이다. \textbf{grounding 주의
(\AL{7})}: 그 기법은 \emph{Fredholm} integral equation of the second kind 를 대상으로
하므로, 본 형식이 \emph{Volterra} 인지 Fredholm 인지에 따른 적용성을 검증 대상(DQ-v3-2)
으로 둔다.

\subsection{Grounded spine 요약 (S1--S14, synthesis)}\label{sec:spine}
\begin{longtable}{p{0.06\textwidth} p{0.66\textwidth} p{0.2\textwidth}}
\toprule ID & 골격 & 상태 [AL] \\ \midrule \endhead
S1 & charge conservation $\Rightarrow$ implicit $V_n$ & GROUNDED \AL{9} \\
S2 & equilibrium target $\theta_\eq$ (일반형, smooth; lattice-gas 특수해) & GROUNDED \AL{3a} \\
S4 & activation free energy $\Delta G_a(T)$ & GROUNDED \AL{1} \\
S5 & ★ effective barrier $\Delta G_\eff=\Delta G_a-\chi\mathcal A$ & BOUNDED \AL{2} \\
S6 & single-mode relaxation $\dot\theta=k(\theta_\eq-\theta)$ & BOUNDED \AL{5} \\
S7 & single-mode exponential kernel $L=v_Q/k$ & GROUNDED \AL{1} \\
S8 & ★ barrier dist.\ $\rho_G\to$ relaxation-length spectrum $A_L$ & BOUNDED \AL{6} \\
S9 & ★ kernel integral for observed tail & BOUNDED \AL{6} \\
S10 & ★ Refs 6/7 propagator = spectrum integral closure & GROUNDED+DQ \AL{7} \\
S11 & ICA mapping $dQ/dV_n$ & GROUNDED \AL{9} \\
S12 & ★ spectrum 의 $T/\psi$ shift = stretched tail & FLAGGED novel \AL{8} \\
S13 & ★ falsification + non-uniqueness 처리 & method \AL{10} \\
S14 & fitting model + identifiability + validation & GROUNDED \AL{10} \\
\bottomrule
\end{longtable}
\textbf{AGP.} 모든 assumption 은 AL-\# 인용. FLAGGED 는 established 로 사용 금지; BOUNDED
는 유효범위 병기. step-function·$0\!\to\!1$ 급점프·$\max$·$\min$·Heaviside·hard-switch
정의식 금지 (smooth only).

% ====================================================================
\section{Notation 과 convention}\label{sec:notation}
% ====================================================================
\begin{longtable}{p{0.16\textwidth} p{0.14\textwidth} p{0.62\textwidth}}
\toprule 기호 & 단위 & 의미 \\ \midrule \endhead
$q$ & --- & 정규화 진행 좌표 $q=Q_\ext/Q_\cell$ \\
$Q_\ext$ & C & 외부 누적 charge; $v_Q\equiv dQ_\ext/dt=|I|>0$ (scan speed) \\
$Q_\cell$ & C & 기준 cell capacity \\
$T,F,R,k_B,h$ & --- & temperature, Faraday, gas const, Boltzmann, Planck \\
$V_n$ & V & internal anode (analysis) potential; charge conservation 의 root \\
$V_{n,\app}$ & V & apparent(measured) potential $=V_n+s_I|I|R_n$ \\
$V_{n,\dimedrive}$ & V & driving potential (rate/barrier 인수); 기본 $=V_n$, reduced $\approx V_{n,\app}$ \\
$R_n(q,T,|I|)$ & $\Omega$ & 유효 polarization 계수 \\
$j,\,N_p$ & --- & effective transition index, 개수 \\
$\theta_j,\ \theta_{\eq,j}$ & --- & local mode 진행률, equilibrium target ($0\!\le\!\theta\!\le\!1$) \\
$r_j$ & --- & kinetic lag $=\theta_{\eq,j}-\theta_j$ \\
$U_j(T),w_j(T)$ & V & equilibrium 중심 전위, isotherm width \\
$\Delta H_{a},\Delta S_{a},\Delta G_a$ & J/mol & activation enthalpy/entropy/free energy \\
$\mathcal A_j$ & J/mol & driving force $=s_{\phi,j}F(V_{n,\dimedrive}-U_j)$ \\
$\chi_j$ & --- & transfer coefficient (BEP), $0\le\chi_j\le1$ \\
$W_\psi$ & J/mol & potential barrier-lowering work $=\chi_j\mathcal A_j$ \\
$\Delta G_{\eff}$ & J/mol & ★ effective barrier $=\Delta G_a-\chi\mathcal A=G-W_\psi$ \\
$G$ & J/mol & local activation barrier (mode 별; hidden rate variable) \\
$\nu(T),\kappa(T),k_0(T)$ & 1/s,--,1/s & Eyring prefactor $\nu=(k_BT/h)\kappa$ \\
$k$ & 1/s & rate constant \\
$L$ & --- & local relaxation length (in $q$) $=v_Q/k$ (단, $q$ 정규화) \\
$\rho_G(G;T,\psi)$ & 1/(J/mol) & activation barrier distribution \\
$A_L(L;T,\psi)$ & 1/L & relaxation-length spectral \emph{density} (in $L$; $\int A_L\,dL$ 가 weight) \\
$\Theta$ & --- & 전체 effective phase progress; $Q_p$ = 그 capacity scale \\
$C_\bg(V_n,T)$ & C/V & background chemical capacitance \\
\bottomrule
\end{longtable}
\begin{groundingbox}
\textbf{Smooth-only.} 모든 정의식은 smooth. 금지 모델의 대표 예: $\theta=H(G_c-G)$
(Heaviside threshold completion) — barrier 를 넘으면 state 가 급점프한다는 모델. 본 chapter
는 이를 채택하지 않는다 (사용자 2026-05-28; AGP-3). 단위: $\Delta G$ J/mol, $\Delta G/RT$
무차원, $k,\nu$ 1/s, $V$ V.
\end{groundingbox}

% ====================================================================
\section{Graphite staging 과 effective transition}\label{sec:staging}
% ====================================================================
Graphite lithiation 은 stage 4 $\to$ 3 $\to$ 2L $\to$ 2 $\to$ 1 sequence 로 진행하며 각
transition 이 $dQ/dV_n$ 에 peak 를 남긴다 \cite{dahn1991,ohzuku1993,funabiki1999} (\AL{4}).
인접 stage 의 중심 전위 차가 width 이하이면 fusion 되어 observable peak 수가 $N_p\approx
3\sim4$ 가 된다. 본 Chapter 는 이를 effective transition $j=1,\dots,N_p$ 로 둔다. 단 각
effective transition 내부에도 particle/domain/boundary 다양성에 의한 \emph{local mode 분포}
가 존재하며, 이것이 \S\ref{sec:spectrum} 의 spectrum 의 물리적 근원이다.

% ====================================================================
\section{Charge conservation 과 implicit $V_n$}\label{sec:charge}
% ====================================================================
\textbf{물리.} anode 내부 charge = 외부 charge (conservation). 이것이 $V_n$ 을 결정한다
(외부 함수 아님; \AL{9}).

\textbf{Derivation D1.} background 와 phase progress 기여의 합
\begin{equation}
Q_\cell\,q = Q_\bg(V_n,T) + Q_p\,\Theta(V_n,q,T),\qquad \Theta=\frac1{N_p}\sum_j w^\Theta_j\theta_j,
\label{eq:charge_balance}
\end{equation}
를 $V_n$ 에 대해 풀어 $V_n$ 을 얻는다 (single-particle/porous-electrode 의 implicit $\mu$
관계와 동형 \cite{doyle1993,newman}). unique smooth 해 조건:
\begin{equation}
\frac{\partial Q_\bg}{\partial V_n}\ge\epsilon_Q>0.
\label{eq:monotone}
\end{equation}
$\theta_j=\theta_{\eq,j}$ 이면 그 해를 $V_n=V_{n,\mathrm{OCV}}(q,T)$ 로 정의 — 외부 OCV
lookup 은 이 charge-balance 의 equilibrium 특수해다.

% ====================================================================
\section{Equilibrium target $\theta_\eq$ (일반형, smooth)}\label{sec:equilibrium}
% ====================================================================
\textbf{보수적 입장.} 본 Chapter 는 equilibrium target 의 함수 형태를 특정하지 않는다.
요구 조건은 (i) smooth, (ii) $0\le\theta_\eq\le1$, (iii) $V_n$ 의 단조 함수뿐이다:
\begin{equation}
\theta_{\eq,j}=\theta_{\eq,j}(V_n,T),\qquad 0\le\theta_{\eq,j}\le1.
\label{eq:thetae}
\end{equation}
(사용자 정정: ``가우시안''은 sharp peak 예시일 뿐 형태 확정 아님 → 특정 형태 강제 X.)

\textbf{Grounded 특수해 (lattice-gas, \AL{3a}).} 명시적 형태가 필요하면, graphite
intercalation 의 grounded thermodynamics 인 lattice-gas / regular-solution chemical
potential \cite{mckinnon1983,bazant2013} 을 쓴다: $\mu(\theta)=\mu^0+RT\ln[\theta/(1-\theta)]
+\Omega_j(1-2\theta)$, electrochemical equilibrium $\mu=s_{\phi,j}F(V_n-U_j^0)$ 에서
($\Omega_j=0$ ideal limit)
\begin{equation}
\theta_{\eq,j}=\big[1+\exp(-s_{\phi,j}(V_n-U_j)/w_j)\big]^{-1},\qquad w_j=RT/F,
\label{eq:latticegas}
\end{equation}
즉 smooth logistic ($0\!\to\!1$ 급점프 아님). $\Omega_j\ne0$ 이면 width/형태 보정, smooth
유지.
\begin{flagbox}
\textbf{형태는 data 결정 (DQ-v3-1).} erf/Gaussian-cumulative isotherm 은 graphite
intercalation 에 대한 thermodynamic derivation 이 문헌에 없으므로 (organic-semiconductor
Gaussian Disorder Model \cite{baessler1993} 과의 analogy + 경험 peak-fit 수준) FLAGGED
ansatz 로만 둔다 (\AL{3b}). 실제 형태는 저속 OCV peak 의 near-peak lineshape 으로 판별
(\S\ref{sec:falsify}). 본 Chapter 의 tail 결론은 $\theta_\eq$ 의 \emph{구체 형태에
의존하지 않는다} — post-peak 에서 $d\theta_\eq/dq\simeq0$ 만 쓰기 때문 (\S\ref{sec:kernel}).
\end{flagbox}

% ====================================================================
\section{★ Effective barrier $\Delta G_\eff$}\label{sec:barrier}
% ====================================================================
\textbf{Derivation D2 (activation free energy, \AL{1}).} transition-state theory
\cite{eyring1935,eyringchemrev1935} 에서 $\Delta G_a(T)=\Delta H_a-T\Delta S_a$.

\textbf{Derivation D3 (potential lowering, BEP; \AL{2}).} 외부 driving force 가 barrier 를
부분적으로 낮춘다는 것은 reaction-coordinate 상 transition state 위치(Hammond)에 따른
비율 $\chi$ 로 표현된다 (Brønsted--Evans--Polanyi \cite{evanspolanyi1938}; Butler--Volmer
transfer coefficient \cite{bardfaulkner}). driving force $\mathcal A_j=s_{\phi,j}F(V_{n,\dimedrive}
-U_j)$ 에 대해 lowering work $W_\psi=\chi_j\mathcal A_j$, 따라서
\begin{equation}
\boxed{\;\Delta G_{\eff}=\Delta G_a-\chi_j\,\mathcal A_j\;=\;G-W_\psi.\;}
\label{eq:Geff}
\end{equation}
여기서 $G\equiv\Delta G_a$ (mode 의 intrinsic barrier; \S\ref{sec:spectrum} 에서 분포),
$W_\psi=\chi_j\mathcal A_j$ (uniform potential lowering). $\mathcal A_j>0$ 이면 barrier 가
낮아진다 — 사용자 verbatim 의 ``극판 전위에 따른 배리어를 낮추는 효과''.
\begin{boundbox}
\textbf{Marcus bound (\AL{2}).} 선형 lowering 은 small driving force 의 leading-order
근사다. Marcus 이론 \cite{marcus1956} 은 barrier 가 driving force 의 2차 함수임을 보이며
큰 $|\mathcal A_j|$ 에서 곡률/inverted region 으로 선형 근사가 깨진다. 본 선형식은
$V_{n,\dimedrive}\approx U_j$ 인 tail 영역에서 유효; 전역 정당화 아님.
\end{boundbox}
\textbf{$\Delta G_\eff$ 가 작아지는 영역 (step 금지).} $W_\psi$ 가 $G$ 에 근접해 activation
이 더 이상 rate-limiting 이 아니면, 식~\eqref{eq:Geff} 를 $\max(\cdot,0)$/clip 으로 고치지
\emph{않는다}. 그 영역에서는 attempt frequency, site availability, phase-boundary motion,
diffusion, transport, current conservation 같은 \emph{다른 physical bottleneck} 이 rate 를
제한한다 (\S\ref{sec:falsify}).

% ====================================================================
\section{Rate constant 과 single-mode exponential kernel}\label{sec:rate}
% ====================================================================
\textbf{Derivation D4 (Eyring rate, \AL{1}).}
\begin{equation}
k(G,T,\psi)=k_0(T)\exp\!\Big[-\frac{G-W_\psi}{RT}\Big],\qquad k_0(T)=\frac{k_BT}{h}\kappa(T).
\label{eq:rate}
\end{equation}
prefactor 는 $k_BT/h$ (Eyring) 이며 임의 상수가 아니다.

\textbf{Derivation D5 (single-mode relaxation, \AL{5}).} 한 local mode 의 lag $r=\theta_\eq-\theta$.
forward/backward rate $r_+,r_-$ 의 net flux $\dot\theta=r_+(1-\theta)-r_-\theta$ 는 stationary
target $\theta_\eq=r_+/(r_++r_-)$ 근처에서 $\dot\theta=k(\theta_\eq-\theta)$, $k=r_++r_-$ 로
환원된다. 즉 $\theta_\eq$ 는 target, $k$ 는 그 target 으로 접근하는 mobility — 이 구분이
double counting 을 막는다 (target 과 mobility 를 동시에 자유 fitting 하지 않음).
\begin{boundbox}
\textbf{유효범위 (\AL{5}).} 선형 relaxation 은 equilibrium 근처 1차 Taylor 이며 tail 영역
($\theta\approx\theta_\eq$) 에서 일반성을 잃지 않는다 \cite{degrootmazur1962,onsager1931}.
\end{boundbox}

\textbf{Derivation D6 ($q$-coordinate, single-mode kernel).} $v_Q=|I|$, $dq/dt=|I|/Q_\cell$
에서 $d\theta/dq=(Q_\cell k/|I|)\,r$, 그리고
\begin{equation}
\frac{dr}{dq}+\frac{Q_\cell k}{|I|}\,r=\frac{d\theta_\eq}{dq}.
\label{eq:r_ode}
\end{equation}
post-peak 에서 $d\theta_\eq/dq\simeq0$ 이면
\begin{equation}
\boxed{\;r(q)\simeq r(q_a)\exp\!\Big[-\frac{q-q_a}{L}\Big],\qquad L=\frac{|I|}{Q_\cell\,k}.\;}
\label{eq:single_kernel}
\end{equation}
이것은 관측 tail 전체가 단일 exponential 이라는 뜻이 \emph{아니다} — \textbf{하나의 local
mode 가 만드는 exponential kernel} 이다. 차원: $L$ 은 무차원 ($q$ 단위). $|I|/(Q_\cell k)=$
A/(C$\cdot$1/s)$=$무차원. 일치.

% ====================================================================
\section{★ Barrier distribution $\to$ relaxation-length spectrum}\label{sec:spectrum}
% ====================================================================
\textbf{물리.} 실제 electrode 에는 단일 $G$ 가 아니라, particle size·local staging
environment·domain boundary·defect·stress·surface·transport 차이로 인한 \emph{barrier
distribution} $\rho_G(G;T,\psi)$ 가 있다. 그러나 관측 tail 의 shape 은 $\rho_G$ 와 같지
않다 — $G$ 가 먼저 rate 로, 다시 length 로 변환되기 때문이다.

\textbf{Derivation D7 (barrier$\to$length 지수 매핑).} 식~\eqref{eq:rate} 와 $L=|I|/(Q_\cell k)$
에서
\begin{equation}
\boxed{\;L(G,T,\psi)=\frac{|I|}{Q_\cell\,k_0(T)}\exp\!\Big[\frac{G-W_\psi}{RT}\Big].\;}
\label{eq:L_of_G}
\end{equation}
$G$ 의 작은 변화가 $L$ 에서 \emph{지수적으로 확대}된다 — 이것이 stretched tail 의 핵심
근원이다 (long tail 은 $\rho_G$ 자체가 길어서가 아니라 barrier$\to$length 매핑이 지수이기
때문일 수 있다). 역변환 $G(L)=W_\psi+RT\ln[Q_\cell k_0 L/|I|]$, Jacobian $|dG/dL|=RT/L$.
따라서 relaxation-length spectrum:
\begin{equation}
A_L(L;T,\psi)=\rho_G\big(G(L,T,\psi);T,\psi\big)\,\frac{RT}{L}\,A_0(L;T,\psi),
\label{eq:spectrum}
\end{equation}
$A_0$ 는 각 mode 의 initial residual amplitude·population weight·accessibility (무차원
per-mode factor). $A_L$ 의 $1/L$ 차원은 $RT/L$ Jacobian 에서 오며, $A_L$ 은 weight 가
아니라 $L$ 에 대한 \emph{density} 다 ($\int A_L\,dL=$ 총 weight).
\begin{boundbox}
\textbf{Non-uniqueness (\AL{6}).} barrier distribution $\to$ stretched kinetics 의 역매핑은
\emph{유일하지 않다} (Plonka \cite{plonka2001}; retrapping 등 다른 메커니즘도 동일 stretched
tail 재현). 따라서 본 Chapter 는 $\rho_G$ 를 관측 tail 에서 \emph{역산하지 않고},
주어진 (또는 독립 측정한) $\rho_G$ 로부터 tail 을 \emph{forward 예측}하는 데에만 쓴다
(\S\ref{sec:falsify}).
\end{boundbox}

% ====================================================================
\section{★ Kernel integral for the observed tail}\label{sec:kernel}
% ====================================================================
전체 phase progress 의 post-peak tail 기여는 local exponential kernel 의 mixture:
\begin{equation}
\boxed{\;\frac{d\Theta_\mathrm{tail}}{dq}\simeq\int_0^\infty A_L(L;T,\psi)\,\frac1L\,
\exp\!\Big[-\frac{q-q_a}{L}\Big]\,dL.\;}
\label{eq:kernel_integral}
\end{equation}
이것이 본 Chapter 의 중심식이다 (Codex 산출물과의 대조에서 채택한 spectrum 관점). single-
mode 식~\eqref{eq:single_kernel} 은 kernel 하나일 뿐이며, 관측 stretched tail 은 spectrum
$A_L$ 의 large-$L$ 영역이 결정한다. 식~\eqref{eq:kernel_integral} 은 Laplace-type mixture
이므로 $\rho_G$ 의 형태에 따라 단일 지수~stretched~power-law 사이를 연속적으로 잇는다
(smooth). 일반적인 두 시점 kernel 로 쓰면
\begin{equation}
\mathcal K(q,q')=\int_0^\infty A_L(L;T,\psi)\,\frac1L\,\exp\!\Big[-\frac{q-q'}{L}\Big]\,dL,
\label{eq:K_def}
\end{equation}
이고 식~\eqref{eq:kernel_integral} 은 $\mathcal K(q,q_a)$ 에 해당한다 (이 $\mathcal K$ 가
\S\ref{sec:closure} 의 closure 에 들어가는 kernel).

\textbf{상태 주의 (\AL{6}).} 본 spectrum/kernel 관점은 Assumption Ledger 에서 BOUNDED
이며, 식~\eqref{eq:kernel_integral} 을 본 Chapter 의 중심으로 쓰는 것은 \emph{spectrum
hypothesis 가 \S\ref{sec:falsify} 의 N3/N4 falsification 을 통과한다는 전제} 하에서다
(single-mode 식~\eqref{eq:single_kernel} 은 $A_L=\delta(L-L_0)$ 특수해로 항상 포함).

% ====================================================================
\section{★ Refs 6/7 propagator: spectrum integral 의 closure}\label{sec:closure}
% ====================================================================
\textbf{왜 필요한가.} 식~\eqref{eq:kernel_integral} 은 $A_L$ 을 알면 forward 평가 가능
하지만, $\theta_j$ 와 $V_n$ 이 charge balance (식~\eqref{eq:charge_balance}) 로 결합되어
있어 일반적으로 $\Theta$ 가 적분 안팎에 동시에 나타나는 \emph{self-consistent integral
equation} 이 된다. 사용자 박사 연구 \cite{lee2011,son2013} 의 propagator / ratio-
substitution 이 이를 닫는다.

\textbf{Derivation D8 (ratio-substitution closure).} 시간 적분형
$\Theta(q)=\Theta_0+\int_0^q \mathcal K(q,q')[\Theta_\eq(q')-\Theta(q')]\,dq'$ 에서 미지
$\Theta(q')$ 에 비율 ansatz $\Theta(q')/\Theta(q)\approx\Theta^{\simplecase}(q')/\Theta^{\simplecase}(q)$
($\Theta^{\simplecase}$ = single reference mode 의 해, $L$ 고정) 를 적용하면 $\Theta(q)$ 가
적분 밖으로 나와
\begin{equation}
\boxed{\;\Theta(q)=\frac{\Theta_0+\int_0^q \mathcal K\,\Theta_\eq\,dq'}
{1+\dfrac{1}{\Theta^{\simplecase}(q)}\int_0^q \mathcal K\,\Theta^{\simplecase}\,dq'}.\;}
\label{eq:closed}
\end{equation}
$\mathcal K(q,q')$ 는 식~\eqref{eq:K_def} 의 spectrum kernel (식~\eqref{eq:kernel_integral}
이 $\mathcal K(q,q_a)$).
\begin{boundbox}
\textbf{적용성 (\AL{7}, DQ-v3-2).} Refs 6/7 의 기법은 \emph{Fredholm} 2nd-kind 를 대상으로
한다. 본 형식이 시간 상한 $q$ 의 \emph{Volterra} 인 한, 식~\eqref{eq:closed} 는 근사·검증-
tier 결과이며 \emph{direct numerical integration of} 식~\eqref{eq:kernel_integral} \emph{대비
오차로 검증}해야 한다 (가정으로 확정하지 않음). ratio ansatz 는 $\Theta_0=0$ reference 에서
self-consistent.
\end{boundbox}

% ====================================================================
\section{ICA mapping}\label{sec:ica}
% ====================================================================
background storage $dQ_\bg=C_\bg(V_n,T)\,dV_n$ 와 phase progress $Q_p\,d\Theta$ 의 합
$dQ=C_\bg dV_n+Q_p d\Theta$ 에서
\begin{equation}
\boxed{\;\frac{dQ}{dV_n}=\frac{C_\bg(V_n,T)}{1-Q_p\,(d\Theta/dQ)}.\;}
\label{eq:ica}
\end{equation}
식~\eqref{eq:kernel_integral} 의 tail 기여가 $d\Theta/dQ$ 에 들어가 ICA tail 로 나타난다.
관측은 $V_{n,\app}$ 축이므로 voltage-axis tail scale 은 $L_\varphi\simeq|dV_n/dQ|_{q_a}L$
로 local DVA mapping 에 의해 변형된다. $V_{n,\app}=V_n+s_I|I|R_n$, $V_n$ = charge-balance
root (식~\eqref{eq:charge_balance}); rate/barrier 인수 $V_{n,\dimedrive}=V_n$ (기본).

% ====================================================================
\section{★ Spectrum 의 $T/\psi$ 의존 = stretched tail (O2)}\label{sec:tailT}
% ====================================================================
\textbf{Derivation D9 (novel; \AL{8}).} 식~\eqref{eq:L_of_G} 에서:
\begin{itemize}
\item \textbf{Low $T$}: $\exp[(G-W_\psi)/RT]$ 가 커지고 $k_0(T)$ 가 작아져 같은 $G-W_\psi$
  에 대해 $L$ 이 커진다 $\Rightarrow$ $A_L$ weight 가 large-$L$ 로 이동 $\Rightarrow$
  식~\eqref{eq:kernel_integral} 의 long-$q$ tail 이 길어진다 (저온 긴 tail).
\item \textbf{High $T$}: $RT$ 증가로 $L$ 이 작아져 $A_L$ 이 short-$L$ 로 이동 $\Rightarrow$
  tail 이 짧게 끝난다.
\item \textbf{Potential assistance}: $W_\psi=\chi_j\mathcal A_j$ 증가 시 $L$ 감소,
  \begin{equation}
  \frac{\partial\ln L}{\partial V_{n,\dimedrive}}=-\frac{\chi_j s_{\phi,j}F}{RT}<0\ (\text{forward}),
  \label{eq:psi_shift}
  \end{equation}
  즉 potential 이 spectrum 을 short-$L$ 로 shift $\Rightarrow$ tail 단축. 이는 state jump
  가 아니라 \emph{spectrum shift} 다.
\end{itemize}
\textbf{Equilibrium-width 와의 구별.} homogeneous lattice-gas width 는 $w_j=RT/F$ 라 저온서
오히려 좁아진다. 따라서 ``저온 긴 tail''은 homogeneous-equilibrium 효과로 설명되지 않으며
\emph{kinetic spectrum shift} 가 필요하다 (단, heterogeneous disorder width 는 $T$-약의존
이므로 별도 — \S\ref{sec:falsify} lineshape test). 본 절은 kinetic-spectrum 메커니즘이
(O2) 와 \emph{강하게 시사(consistent)} 됨을 보이며, \emph{확정}은 falsification 통과 전제.

% ====================================================================
\section{★ Falsification 과 non-uniqueness 처리}\label{sec:falsify}
% ====================================================================
\textbf{왜 필요한가.} 관측 tail 을 barrier-spectrum 에 귀속하려면 competing source 를
배제해야 한다 (\AL{10}; ICA 가 rate 로 broadening 된다는 것은 알려져 있으므로
\cite{dubarry2022,fly2020} 단순 rate-broadening 과 구별이 핵심).
\begin{longtable}{p{0.24\textwidth} p{0.40\textwidth} p{0.30\textwidth}}
\toprule Competing source & tail 기여 & 판별 \\ \midrule \endhead
Equilibrium width & homogeneous/disorder broadening & 저속 OCV peak near-peak lineshape; 저온서 좁아지면 kinetic 아님 (homogeneous 한정) \\
Polarization $R_n$ & $V_{n,\app}$ shift & pulse/current-interrupt/다중 $|I|$ \\
Transport/diffusion & 저온 느린 확산 lag & GITT relaxation signature; $D(T)$ Arrhenius vs barrier slope \\
Barrier distribution $\rho_G$ \cite{plonka2001} & spectrum kernel & \emph{비유일} 역매핑 — forward prediction 만; 역산 금지 \\
\bottomrule
\end{longtable}

\textbf{★ 비퇴화 discriminator.} tail length scaling 만으로는 부족하다: $L\propto|I|$ 는
transport lag 도 공유한다 (퇴화). barrier-spectrum 의 고유 signature 는 식~\eqref{eq:L_of_G},
\eqref{eq:psi_shift} 의 \emph{potential-dependent} 항이다. 즉 (a) tail length 의 Arrhenius
slope 자체가 $(V_{n,\dimedrive}-U_j)$ 에 의존하고, (b) potential/$|I|$ 증가 시 spectrum 이
$\chi_j$ (와 $R_n$) 를 통해 short-$L$ 로 shift 한다. transport/단순 polarization 은 이런
potential-dependent slope 를 갖지 않는다.

\textbf{Null-result 규칙.}
\begin{enumerate}[label=(N\arabic*)]
\item tail length 의 $\ln$ vs $1/T$ slope 가 독립 측정 $E_D$ 와 일치하고 potential 의존이
  없으면 $\Rightarrow$ barrier-spectrum 귀속 기각 (transport-dominant).
\item potential/$|I|$ 증가 시 spectrum/peak shift 가 없으면 $\Rightarrow$ $\chi_j$ 항 부재
  $\Rightarrow$ potential-lowering 가설 기각.
\item 저속 OCV near-peak lineshape 이 저온서 좁아지지 않으면 $\Rightarrow$ heterogeneous
  disorder 기여 존재 — kinetic 단독 귀속 기각, disorder 와 분리.
\item $\rho_G$ 를 관측 tail 에서 역산하지 않는다 (non-unique); $\rho_G$ 는 독립 근거
  (입자분포·EIS·계산) 로 주고 tail 을 forward 예측해 비교만 한다.
\end{enumerate}

% ====================================================================
\section{Fitting 모델, identifiability, validation}\label{sec:fitting}
% ====================================================================
\textbf{Deliverable.} ICA tail 의 closed-form 논리식 (평가 순서 inner$\to$outer):
$k(G,T,\psi)$ (식~\eqref{eq:rate}) $\to$ $L(G)$ (식~\eqref{eq:L_of_G}) $\to$ spectrum
$A_L$ (식~\eqref{eq:spectrum}) $\to$ kernel integral $d\Theta_\mathrm{tail}/dq$
(식~\eqref{eq:kernel_integral}, Refs 6/7 closure 식~\eqref{eq:closed}) $\to$ $dQ/dV_n$
(식~\eqref{eq:ica}). \textbf{유효범위}: $\Delta G_\eff\ge0$ (Marcus bound), tail 영역
$V_{n,\dimedrive}\approx U_j+O(w_j)$.

\textbf{Parameter / 실험 매핑 (\AL{10}).}
\begin{longtable}{p{0.26\textwidth} p{0.30\textwidth} p{0.38\textwidth}}
\toprule parameter & 의미 & 결정 실험 \\ \midrule \endhead
$U_j(T),w_j(T),Q_p$ & equilibrium target & 저속 OCV / GITT 준평형 \\
$\Delta H_a,\Delta S_a$ & activation & 다중 $T$ Arrhenius \\
$\chi_j$ (혹은 $\Lambda_\psi$) & potential lowering & 다중 $|I|$/potential shift \\
$\rho_G$ (외부 입력) & barrier spectrum & 독립 근거(입자분포/EIS/계산)로 \emph{주어짐} — tail 역산 금지 (forward only) \\
$k_0(T)$ & prefactor & Arrhenius intercept \\
\bottomrule
\end{longtable}
\textbf{Non-identifiability 경고.} $R_n,\rho_G,w_j,k_0$ 는 모두 broadening 에 기여하므로
단일 $T$/제한된 $|I|$ 로 동시 분리 불가. 독립 실험(pulse/GITT/EIS, 다중 $T\times|I|$) 으로
제약 필수. fitting solver 구현은 범위 외 (P1).

% ====================================================================
\section{종합·grounding 감사·참고문헌}\label{sec:summary}
% ====================================================================
\subsection{결과 요약}
\begin{enumerate}[label=\textbf{(R\arabic*)}]
\item charge conservation 으로 $V_n$ implicit (S1, \AL{9}).
\item equilibrium target 일반형(smooth); lattice-gas grounded 특수해, 형태 data 결정 (S2, \AL{3a}).
\item ★ effective barrier $\Delta G_\eff=\Delta G_a-\chi\mathcal A=G-W_\psi$, Marcus-bounded (S5, \AL{2}).
\item single-mode relaxation + exponential kernel $L=|I|/(Q_\cell k)$ (S6/S7, \AL{1},\AL{5}).
\item ★ barrier distribution $\to$ 지수 매핑 $\to$ relaxation-length spectrum $A_L$ (S8, \AL{6}).
\item ★ 관측 tail = spectrum kernel integral (S9); Refs 6/7 propagator 로 closure (S10, \AL{7}, 검증-tier).
\item ★ stretched tail 의 $T/\psi$ 의존 = spectrum shift (novel; S12, \AL{8}); homogeneous-equilibrium 으로 설명 불가, kinetic-spectrum 강하게 시사 (확정은 falsification 전제).
\item ★ falsification protocol + 비퇴화 discriminator + $\rho_G$ non-unique 역산 회피 (S13, \AL{10}).
\item fitting 논리식 + identifiability + validation (S14).
\end{enumerate}

\subsection{Grounding 감사 (AGP self-check)}
\begin{longtable}{p{0.6\textwidth} p{0.32\textwidth}}
\toprule 항목 & 상태 \\ \midrule \endhead
모든 본문 assumption 이 AL-\# 인용 & 통과 \\
FLAGGED(AL-3b erf, AL-7 적용성, AL-8 tail) established 미사용 & 통과 \\
BOUNDED(AL-2 Marcus,5 near-eq,6 non-unique) 유효범위 병기 & 통과 \\
step-function/$0\!\to\!1$ 급점프 정의식 부재 & 통과 (smooth only) \\
activation barrier 유지 (사용자 2026-05-28) & 통과 (\S\ref{sec:barrier}) \\
``가우시안''을 형태 확정으로 쓰지 않음 & 통과 (\S\ref{sec:equilibrium}) \\
tail 귀속을 ``확정'' 아닌 ``시사+protocol'' & 통과 (\S\ref{sec:tailT},\ref{sec:falsify}) \\
$\rho_G$ non-unique 역산 회피 (forward only) & 통과 (\S\ref{sec:spectrum},\ref{sec:falsify}) \\
사용자 PhD Refs 6/7 = kernel closure 로 사용 & 통과 (\S\ref{sec:closure}) \\
\bottomrule
\end{longtable}

\subsection{Chapter 2+ 로 전달}
single-mode relaxation, spectrum kernel, $\Delta G_\eff$ 가 reversible/irreversible heat,
$dV/dT$, Butler--Volmer kinetics, 충방전 hysteresis(branch-dependent spectrum), full-cell
의 input (별도 roadmap).

\begin{thebibliography}{99}
\bibitem{eyring1935} H. Eyring, ``The activated complex in chemical reactions,''
  \emph{J. Chem. Phys.} \textbf{3}, 107 (1935). DOI: 10.1063/1.1749604.
\bibitem{eyringchemrev1935} S. Glasstone, K. J. Laidler, H. Eyring, ``The activated
  complex and the absolute rate of chemical reactions,'' \emph{Chem. Rev.} \textbf{17},
  301 (1935). DOI: 10.1021/cr60056a006.
\bibitem{evanspolanyi1938} M. G. Evans, M. Polanyi, ``Inertia and driving force of
  chemical reactions,'' \emph{Trans. Faraday Soc.} \textbf{34}, 11 (1938). DOI:
  10.1039/TF9383400011.
\bibitem{marcus1956} R. A. Marcus, ``On the theory of oxidation--reduction reactions
  involving electron transfer,'' \emph{J. Chem. Phys.} \textbf{24}, 966 (1956). DOI:
  10.1063/1.1742723.
\bibitem{bardfaulkner} A. J. Bard, L. R. Faulkner, \emph{Electrochemical Methods},
  Wiley (Ch.~3).
\bibitem{mckinnon1983} W. R. McKinnon, R. R. Haering, ``Physical mechanisms of
  intercalation,'' \emph{Modern Aspects of Electrochemistry} \textbf{15}, 235 (1983).
\bibitem{bazant2013} M. Z. Bazant, ``Theory of chemical kinetics and charge transfer
  based on nonequilibrium thermodynamics,'' \emph{Acc. Chem. Res.} \textbf{46}, 1144
  (2013). DOI: 10.1021/ar300145c.
\bibitem{dahn1991} J. R. Dahn, ``Phase diagram of Li$_x$C$_6$,'' \emph{Phys. Rev. B}
  \textbf{44}, 9170 (1991). DOI: 10.1103/PhysRevB.44.9170.
\bibitem{ohzuku1993} T. Ohzuku, Y. Iwakoshi, K. Sawai, \emph{J. Electrochem. Soc.}
  \textbf{140}, 2490 (1993). DOI: 10.1149/1.2220849.
\bibitem{funabiki1999} A. Funabiki, M. Inaba, T. Abe, Z. Ogumi, ``Nucleation and
  phase-boundary movement upon stage transformation in lithium--graphite intercalation
  compounds,'' \emph{Electrochim. Acta} \textbf{45}, 865 (1999). DOI:
  10.1016/S0013-4686(99)00290-X.
\bibitem{degrootmazur1962} S. R. de Groot, P. Mazur, \emph{Non-Equilibrium
  Thermodynamics}, North-Holland (1962); Dover (1984), Ch.~X.
\bibitem{onsager1931} L. Onsager, ``Reciprocal relations in irreversible processes,''
  \emph{Phys. Rev.} \textbf{37}, 405 (1931).
\bibitem{plonka2001} A. Plonka, \emph{Dispersive Kinetics}, Springer (2001).
\bibitem{lee2011} S. Lee, C. Y. Son, J. Sung, S. H. Chong, ``Communication: Propagator
  for diffusive dynamics of an interacting molecular pair,'' \emph{J. Chem. Phys.}
  \textbf{134}, 121102 (2011). DOI: 10.1063/1.3565476.
\bibitem{son2013} C. Y. Son, J. Kim, J.-H. Kim, J. S. Kim, S. Lee, ``An accurate
  expression for the rates of diffusion-influenced bimolecular reactions with long-range
  reactivity,'' \emph{J. Chem. Phys.} \textbf{138}, 164123 (2013). DOI: 10.1063/1.4802584.
\bibitem{doyle1993} M. Doyle, T. F. Fuller, J. Newman, ``Modeling of galvanostatic
  charge and discharge,'' \emph{J. Electrochem. Soc.} \textbf{140}, 1526 (1993). DOI:
  10.1149/1.2221597.
\bibitem{newman} J. Newman, K. E. Thomas-Alyea, \emph{Electrochemical Systems}, 3rd ed.,
  Wiley (2004).
\bibitem{baessler1993} H. Bässler, ``Charge transport in disordered organic
  photoconductors,'' \emph{Phys. Status Solidi B} \textbf{175}, 15 (1993). DOI:
  10.1002/pssb.2221750102.
\bibitem{dubarry2022} M. Dubarry, D. Anseán, ``Best practices for incremental capacity
  analysis,'' \emph{Front. Energy Res.} \textbf{10}, 1023555 (2022). DOI:
  10.3389/fenrg.2022.1023555.
\bibitem{fly2020} A. Fly, R. Chen, ``Rate dependency of incremental capacity analysis
  (dQ/dV),'' \emph{J. Energy Storage} \textbf{29}, 101329 (2020). DOI:
  10.1016/j.est.2020.101329.
\end{thebibliography}

\end{document}
